\chapter{Variáveis aleatórias e a distribuição binomial}\label{ch:g12:randvar}

Uma \omterm{def:g12:randvar:rv}{variável aleatória} associa um número a cada resultado de um
experimento: um ganho, uma contagem, uma duração. Sua \omterm{def:g12:randvar:exp}{esperança} é a \omterm{def:g11:stat:mean}{média}
a longo prazo dos valores que ela produz; sua \omterm{def:g12:randvar:exp}{variância} mede a dispersão
desses valores. A estrela deste capítulo é a \omterm{def:g12:randvar:binomial}{distribuição binomial}, que
conta os sucessos em repetições \omterm{def:g12:condprob:indep}{independentes}. As variáveis aleatórias e a
\omterm{def:g12:randvar:binomial}{distribuição binomial} apareceram pela primeira vez nos
\cref{ch:g11:prob,ch:g11:binom}; este capítulo as revisa e as aprofunda,
agora com as ferramentas combinatórias do \cref{ch:g12:comb} à
disposição.

\section{Variáveis aleatórias discretas}

\begin{definition}[Variável aleatória, distribuição]\label{def:g12:randvar:rv}
Uma \emph{variável aleatória}\index{variável aleatória} sobre um \omterm{def:g11:prob:model}{espaço
amostral} finito $\Omega$ é uma \omterm{def:g11:func:function}{função} $X \colon \Omega \to \R$. Sua
\emph{distribuição}\index{distribuição} (ou lei) é o conjunto de seus
valores possíveis $x_1, \dots, x_k$ e das \omterm{def:g10:proba:distribution}{probabilidades}
\[
p_i = \P(X = x_i), \qquad \sum_{i=1}^{k} p_i = 1 .
\]
\end{definition}

\begin{definition}[Esperança, variância, desvio padrão]\label{def:g12:randvar:exp}
A \emph{esperança}\index{esperança} de $X$ é
\[
\E(X) = \sum_{i=1}^{k} p_i\, x_i ,
\]
sua \emph{variância}\index{variância} e seu \emph{\omterm{def:g11:stat:variance}{desvio padrão}} são
\[
\V(X) = \E\bigl((X - \E(X))^2\bigr) = \sum_{i=1}^k p_i\bigl(x_i - \E(X)\bigr)^2,
\qquad
\sigma(X) = \sqrt{\V(X)} .
\]
\end{definition}

\begin{proposition}[Fórmula de König--Huygens]\label{prop:g12:randvar:konig}
\index{fórmula de König--Huygens}
$\V(X) = \E(X^2) - \E(X)^2$.
\end{proposition}

\begin{proof}
Escreva $m = \E(X)$ e desenvolva:
\[
\V(X) = \sum_i p_i (x_i^2 - 2m x_i + m^2)
= \E(X^2) - 2m\sum_i p_i x_i + m^2\sum_i p_i
= \E(X^2) - 2m^2 + m^2 . \qedhere
\]
\end{proof}

\begin{proposition}[Transformação afim]\label{prop:g12:randvar:affine}
Para $a, b \in \R$:
\[
\E(aX + b) = a\,\E(X) + b, \qquad
\V(aX + b) = a^2\,\V(X) .
\]
\end{proposition}

\begin{proof}
A primeira é um rearranjo da soma que define a \omterm{def:g12:randvar:exp}{esperança}. Para a segunda,
$aX + b$ se afasta de sua \omterm{def:g11:stat:mean}{média} por $a(X - \E(X))$, e elevar ao quadrado
multiplica por $a^2$.
\end{proof}

\begin{example}[Jogos equitativos]\label{ex:g12:randvar:game}
Um jogo custa $2$ euros; um dado é lançado, e o jogador recebe o valor
mostrado se este for pelo menos $5$, e nada caso contrário. O ganho $G$
assume os valores $-2$ (\omterm{def:g10:proba:distribution}{probabilidade} $\frac46$), $3$ ($\frac16$), $4$
($\frac16$):
\[
\E(G) = \frac{-8 + 3 + 4}{6} = -\frac{1}{6} < 0 .
\]
Em \omterm{def:g11:stat:mean}{média}, o jogador perde $17$ centavos por partida: o jogo é
desfavorável (como quase todos os jogos reais).
\end{example}

\section{Ensaios de Bernoulli e a distribuição binomial}

\begin{definition}[Distribuição de Bernoulli]\label{def:g12:randvar:bernoulli}
Um \emph{ensaio de Bernoulli}\index{ensaio de Bernoulli} é um experimento
com dois resultados, \emph{sucesso} (\omterm{def:g10:proba:distribution}{probabilidade} $p$) e \emph{fracasso}
($q = 1 - p$). O indicador $X$ do sucesso ($X = 1$ em caso de sucesso, $0$
em caso de fracasso) segue a \emph{\omterm{def:g12:randvar:rv}{distribuição} de Bernoulli}
$\mathcal B(p)$:
\[
\E(X) = p, \qquad \V(X) = p(1-p) .
\]
(De fato, $\E(X) = p$, $\E(X^2) = p$, e König--Huygens dá
$\V(X) = p - p^2$.)
\end{definition}

\begin{definition}[Distribuição binomial]\label{def:g12:randvar:binomial}
Repita um \omterm{def:g12:randvar:bernoulli}{ensaio de Bernoulli} $n$ vezes de maneira independente e seja $X$
o número total de sucessos. A \omterm{def:g12:randvar:rv}{distribuição} de $X$ é a \emph{distribuição
binomial}\index{distribuição binomial} $\mathcal B(n, p)$.
\end{definition}

\begin{omfigure}
\begin{tikzpicture}[font=\small]
  \coordinate (root) at (0,0);
  \node (S) at (1.6,1.6) {S};
  \node (F) at (1.6,-1.6) {F};
  \node (SS) at (3.2,2.4) {S};
  \node (SF) at (3.2,0.8) {F};
  \node (FS) at (3.2,-0.8) {S};
  \node (FF) at (3.2,-2.4) {F};
  \node (SSS) at (4.8,2.8) {S};
  \node (SSF) at (4.8,2.0) {F};
  \node (SFS) at (4.8,1.2) {S};
  \node (SFF) at (4.8,0.4) {F};
  \node (FSS) at (4.8,-0.4) {S};
  \node (FSF) at (4.8,-1.2) {F};
  \node (FFS) at (4.8,-2.0) {S};
  \node (FFF) at (4.8,-2.8) {F};
  \draw[omThm, thick] (root) -- (S);
  \draw[omThm, thick] (S) -- (SS);   \draw[omThm, thick] (SS) -- (SSF);
  \draw[omThm, thick] (S) -- (SF);   \draw[omThm, thick] (SF) -- (SFS);
  \draw[omThm, thick] (root) -- (F);
  \draw[omThm, thick] (F) -- (FS);   \draw[omThm, thick] (FS) -- (FSS);
  \draw[gray] (SS) -- (SSS);
  \draw[gray] (SF) -- (SFF);
  \draw[gray] (FS) -- (FSF);
  \draw[gray] (F) -- (FF);
  \draw[gray] (FF) -- (FFS);
  \draw[gray] (FF) -- (FFF);
  \node[gray, anchor=west] at (5.3,2.8) {$X=3$};
  \node[omThm, anchor=west] at (5.3,2.0) {$X=2$};
  \node[omThm, anchor=west] at (5.3,1.2) {$X=2$};
  \node[gray, anchor=west] at (5.3,0.4) {$X=1$};
  \node[omThm, anchor=west] at (5.3,-0.4) {$X=2$};
  \node[gray, anchor=west] at (5.3,-1.2) {$X=1$};
  \node[gray, anchor=west] at (5.3,-2.0) {$X=1$};
  \node[gray, anchor=west] at (5.3,-2.8) {$X=0$};
\end{tikzpicture}

{\small Três \omterm{def:g12:randvar:bernoulli}{ensaios de Bernoulli}: exatamente $\binom{3}{2} = 3$ dos
$2^3$ caminhos dão dois sucessos (em vermelho), cada um com \omterm{def:g10:proba:distribution}{probabilidade}
$p^2(1-p)$, de modo que $\P(X = 2) = 3p^2(1-p)$.}
\end{omfigure}

\begin{theorem}\label{thm:g12:randvar:binomial}
Se $X \sim \mathcal B(n, p)$, então, para $0 \leq k \leq n$:
\[
\P(X = k) = \binom{n}{k} p^k (1-p)^{n-k},
\]
e
\[
\E(X) = np, \qquad \V(X) = np(1-p) .
\]
\end{theorem}

\begin{proof}
Uma \omterm{def:g12:seq:sequence}{sequência} determinada de resultados com $k$ sucessos e $n-k$ fracassos
tem \omterm{def:g10:proba:distribution}{probabilidade} $p^k(1-p)^{n-k}$ por \omterm{def:g12:condprob:indep}{independência}; o número de tais
\omterm{def:g12:seq:sequence}{sequências} é o número de maneiras de colocar os $k$ sucessos entre os $n$
ensaios, a saber $\binom nk$ (\cref{ch:g12:comb}). Somando sobre as
\omterm{def:g12:seq:sequence}{sequências}, obtemos a fórmula --- e o teorema binomial confirma que
$\sum_k \P(X = k) = (p + q)^n = 1$.

Para a \omterm{def:g12:randvar:exp}{esperança}, usando $k\binom nk = n\binom{n-1}{k-1}$
(\cref{exo:g12:comb:7}):
\[
\E(X) = \sum_{k=1}^{n} k\binom nk p^k q^{n-k}
= np\sum_{k=1}^{n} \binom{n-1}{k-1} p^{k-1} q^{(n-1)-(k-1)}
= np\,(p + q)^{n-1} = np .
\]
A fórmula da \omterm{def:g12:randvar:exp}{variância} se demonstra de modo semelhante com a identidade
$k(k-1)\binom nk = n(n-1)\binom{n-2}{k-2}$, que dá
$\E(X(X-1)) = n(n-1)p^2$, donde
$\V(X) = n(n-1)p^2 + np - (np)^2 = np(1-p)$. (Uma demonstração estrutural
--- a \omterm{def:g12:randvar:exp}{variância} de uma soma de variáveis \omterm{def:g12:condprob:indep}{independentes} --- virá com o
\cref{thm:g12:sums:variance}.)
\end{proof}

\begin{omfigure}
\begin{tikzpicture}
\begin{axis}[
  width=10cm, height=4.8cm,
  ybar, bar width=4pt,
  xlabel={$k$}, ylabel={$\P(X = k)$},
  xmin=-1, xmax=21, ymin=0,
  ytick={0, 0.1, 0.2},
  title={$\mathcal B(20,\ 0.3)$}]
  \addplot[fill=omDef!40, draw=omDef] coordinates {
    (0,0.000798) (1,0.006839) (2,0.027846) (3,0.071604)
    (4,0.130421) (5,0.178863) (6,0.191639) (7,0.164262)
    (8,0.114397) (9,0.065370) (10,0.030817) (11,0.012006)
    (12,0.003859) (13,0.001018) (14,0.000218) (15,0.000037)
    (16,0.000005) (17,0.000001) (18,0) (19,0) (20,0)};
\end{axis}
\end{tikzpicture}

{\small A \omterm{def:g12:randvar:rv}{distribuição} $\mathcal B(20, 0.3)$: \omterm{def:g11:stat:mean}{média} $np = 6$, \omterm{def:g11:stat:variance}{desvio
padrão} $\sqrt{npq} \approx 2.05$.}
\end{omfigure}

\begin{method}[Reconhecer uma situação binomial]\label{met:g12:randvar:recognize}
Verifique os três ingredientes antes de escrever $X \sim \mathcal B(n,p)$:
um \emph{número fixo} $n$ de ensaios; \emph{dois resultados} por ensaio
com a \emph{mesma} \omterm{def:g10:proba:distribution}{probabilidade} de sucesso $p$; \emph{\omterm{def:g12:condprob:indep}{independência}} dos
ensaios (amostragem \emph{com} reposição, ou em uma população grande).
Depois use
\[
\P(X \geq 1) = 1 - (1-p)^n
\]
para ``ao menos um sucesso'', e uma calculadora ou tabelas acumuladas para
$\P(X \leq k)$ em geral.
\end{method}

\begin{example}\label{ex:g12:randvar:atleastone}
Quantas vezes é preciso lançar um dado para ter ao menos $99\%$ de chance
de tirar um seis? Com $X \sim \mathcal B\left(n, \frac16\right)$:
$\P(X \geq 1) = 1 - \left(\frac56\right)^n \geq 0.99$ significa
$\left(\frac56\right)^n \leq 0.01$, \emph{isto é},
$n \geq \frac{\ln 0.01}{\ln(5/6)} \approx 25.3$: a partir de $n = 26$
lançamentos.
\end{example}

\section{Exercícios}

\begin{exercise}[$\star$]\label{exo:g12:randvar:1}
Uma \omterm{def:g12:randvar:rv}{variável aleatória} $X$ assume os valores $-1, 0, 2, 5$ com
\omterm{def:g10:proba:distribution}{probabilidades} $0.3, 0.2, 0.4, 0.1$. Calcule $\E(X)$, $\V(X)$ e
$\sigma(X)$.
\end{exercise}

\begin{exercise}[$\star$]\label{exo:g12:randvar:2}
Uma prova de múltipla escolha tem $10$ questões, cada uma com $4$
alternativas, uma das quais é correta. Um estudante responde
uniformemente ao acaso, de maneira independente. Seja $X$ o número de
respostas corretas.
\begin{enumerate}
  \item Dê a \omterm{def:g12:randvar:rv}{distribuição} de $X$, $\E(X)$ e $\sigma(X)$.
  \item Calcule $\P(X = 0)$, $\P(X = 5)$ e $\P(X \geq 1)$.
\end{enumerate}
\end{exercise}

\begin{exercise}[$\star$]\label{exo:g12:randvar:3}
Uma seguradora cobre $n = 400$ clientes; cada um abre um sinistro durante
o ano com \omterm{def:g10:proba:distribution}{probabilidade} $p = 0.05$, de maneira independente. Seja $X$ o
número de sinistros. Identifique a \omterm{def:g12:randvar:rv}{distribuição} de $X$ e calcule sua
\omterm{def:g12:randvar:exp}{esperança} e seu \omterm{def:g11:stat:variance}{desvio padrão}.
\end{exercise}

\begin{exercise}[$\star\star$]\label{exo:g12:randvar:4}
No jogo do \cref{ex:g12:randvar:game}, o organizador quer um jogo
equitativo ($\E(G) = 0$) mudando o preço de entrada $c$. Determine $c$.
Calcule a \omterm{def:g12:randvar:exp}{variância} do ganho nessa versão equitativa; ``equitativo'' é o
mesmo que ``sem risco''?
\end{exercise}

\begin{exercise}[$\star\star$]\label{exo:g12:randvar:5}
Uma jogadora de basquete converte lances livres com \omterm{def:g10:proba:distribution}{probabilidade} $0.7$.
Ela arremessa $8$ vezes (arremessos \omterm{def:g12:condprob:indep}{independentes}). Calcule a
\omterm{def:g10:proba:distribution}{probabilidade} de que converta: exatamente $6$; ao menos $6$; ao menos um.
Qual é o número mais provável de conversões?
\end{exercise}

\begin{exercise}[$\star\star$]\label{exo:g12:randvar:6}
Uma companhia aérea sabe que cada passageiro com reserva comparece com
\omterm{def:g10:proba:distribution}{probabilidade} $0.9$, de maneira independente. Um voo tem $100$ assentos e
a companhia vende $104$ passagens. Exprima, usando uma \omterm{def:g12:randvar:binomial}{distribuição
binomial}, a \omterm{def:g10:proba:distribution}{probabilidade} de que compareçam mais passageiros do que
assentos, e limite-a numericamente com uma calculadora (dê a expressão
exata).
\end{exercise}

\begin{exercise}[$\star\star$]\label{exo:g12:randvar:7}
Seja $X \sim \mathcal B(n, p)$. Mostre que
\[
\frac{\P(X = k+1)}{\P(X = k)} = \frac{n-k}{k+1}\cdot\frac{p}{1-p},
\]
e deduza que a \omterm{def:g12:randvar:rv}{distribuição} cresce até
$k^* = \floor{(n+1)p}$ e decresce depois (a \emph{moda} da binomial).
\end{exercise}

\begin{exercise}[$\star\star\star$]\label{exo:g12:randvar:8}
\emph{(São Petersburgo, domada.)} Uma moeda honesta é lançada até sair
cara, mas no máximo $10$ vezes. Seja $N$ o número de lançamentos
utilizados; o jogador recebe $2^N$ euros se saiu cara, e $0$ caso
contrário.
\begin{enumerate}
  \item Dê a \omterm{def:g12:randvar:rv}{distribuição} de $N$ restrita aos resultados vencedores:
        $\P(\text{primeira cara no lançamento } k) = 2^{-k}$ para
        $1 \leq k \leq 10$, e verifique a \omterm{def:g10:proba:distribution}{probabilidade} total de ganhar.
  \item Calcule o ganho esperado. O que ele se tornaria sem o limite de
        $10$ lançamentos?
\end{enumerate}
\end{exercise}

\section{Problema: o voo com excesso de reservas}

\begin{problem}[{Problema de fim de semana --- as companhias aéreas
vendem mais assentos do que têm, as fábricas aceitam lotes que mal
inspecionaram, e a distribuição binomial arbitra os dois casos}]
\label{pb:g12:randvar:1}
Uma companhia aérea com $100$ assentos vende alegremente $105$
passagens: cerca de $10\,\%$ dos passageiros nunca comparecem, e
assentos vazios não rendem nada. Até onde a companhia pode
esticar a corda antes que os passageiros preteridos devorem o
lucro? A \omterm{def:g12:randvar:binomial}{distribuição binomial}
(\cref{thm:g12:randvar:binomial}) responde na casa decimal --- e
a mesma maquinaria inspeciona lotes de fábrica, precifica rifas
e mantém as seguradoras solventes. Decisões sob repetição: este
é o trabalho diário da binomial.

\medskip
\noindent\textbf{Parte I --- Fluência.}
\begin{enumerate}
  \item Dê a tabela completa da \omterm{def:g12:randvar:rv}{distribuição} do número $X$ de
        seis em três lançamentos de um dado.
  \item Calcule $\E(X)$ e $V(X)$
        (\cref{def:g12:randvar:exp}).
  \item Uma barraca cobra $1$ euro por três lançamentos e paga
        $2$ euros por seis obtido. Calcule o ganho esperado:
        jogo equitativo?
  \item Treino de conferência (\cref{met:g12:randvar:recognize}):
        o número de cartas de copas em $5$ cartas retiradas
        \emph{sem} reposição é binomial? E com reposição?
        Justifique.
  \item Para $X \sim \mathcal B(20,\ 0.3)$: dê $\E(X)$ e
        $\sigma(X)$, e depois $\E(S)$ e $\sigma(S)$ para a
        pontuação $S = 5X - 10$
        (\cref{prop:g12:randvar:affine}).
\end{enumerate}

\medskip
\noindent\textbf{Parte II --- O voo com excesso de reservas.}
Assentos: $100$. Passagens vendidas: $105$. Cada portador de
passagem comparece com \omterm{def:g10:proba:distribution}{probabilidade} $0.9$, de maneira
independente; seja $X$ o número dos que comparecem.
\begin{enumerate}[resume]
  \item Modelo: justifique $X \sim \mathcal B(105,\ 0.9)$ com a
        lista de verificação --- e confesse o ponto mais frágil
        do modelo (as ausências são mesmo \omterm{def:g12:condprob:indep}{independentes}? pense
        em grupos e em tempestades).
  \item Calcule $\E(X)$ e $\sigma(X)$.
  \item Há passageiros preteridos quando $X \geq 101$: escreva
        $\P(X \geq 101)$ como uma soma explícita de cinco
        termos binomiais.
  \item Avalie a soma (calculadora): que fração dos voos vê ao
        menos um passageiro preterido?
  \item Dinheiro: as cinco passagens extras rendem
        $5 \times 200 = 1\,000$ euros; cada passageiro preterido
        custa $800$ euros de indenização. Calcule o número
        esperado de passageiros preteridos,
        $\sum_{k=101}^{105} (k - 100)\,\P(X = k)$, a indenização
        esperada e o veredicto sobre a política.
  \item Projeto: o regulador tolera
        $\P(\text{preterição}) < 5\,\%$. Testando
        $n = 105, 106, 107, \dots$ passagens, encontre o maior
        $n$ permitido.
  \item O atalho do sino: para $n = 105$, calcule o \omterm{pb:g11:stat:1}{escore z}
        do limiar de preterição,
        $z = \frac{100.5 - \E(X)}{\sigma(X)}$, e compare a
        estimativa grosseira ``cerca de $2\sigma$, logo uns
        $2$--$3\,\%$ na cauda superior'' com sua resposta
        exata. (A curva por trás desse atalho é a estrela dos
        próximos capítulos.)
\end{enumerate}

\medskip
\noindent\textbf{Parte III --- O portão da fábrica.}
Um lote de peças é aceito se uma \omterm{def:g10:proba:sample}{amostra} de $20$ contiver no
máximo uma peça defeituosa.
\begin{enumerate}[resume]
  \item Se a taxa real de defeitos é $2\,\%$ (um lote honesto),
        calcule a \omterm{def:g10:proba:distribution}{probabilidade} de aceitação.
  \item Se a taxa é $10\,\%$ (um lote ruim), calcule-a de novo.
  \item Nomeie os dois riscos do procedimento (o lote honesto
        rejeitado; o lote ruim aceito), leia seus valores nas
        questões 13--14 e diga que mudança única melhora os
        dois ao mesmo tempo --- e a que custo.
  \item Por trás da melhoria: mostre que a \emph{\omterm{def:g10:stats:series}{frequência}}
        observada de defeitos $\frac Xn$ em uma \omterm{def:g10:proba:sample}{amostra} de
        tamanho $n$ tem \omterm{def:g11:stat:variance}{desvio padrão}
        $\sqrt{\frac{p(1-p)}{n}}$ e conclua como a precisão
        varia com o tamanho da \omterm{def:g10:proba:sample}{amostra} (um velho conhecido: o
        $\frac{1}{\sqrt n}$ do ano 10).
\end{enumerate}

\medskip
\noindent\textbf{Parte IV --- Jogos, rifas, reservas.}
\begin{enumerate}[resume]
  \item O jogo de São Petersburgo com limite do
        \cref{exo:g12:randvar:8} tem ganho esperado de $10$
        euros. Contraste em um parágrafo com o paradoxo sem
        limite encontrado no \cref{pb:g11:prob:1}: o que
        exatamente o limite doma, e que preço seria agora
        equitativo?
  \item Uma rifa beneficente vende $200$ bilhetes a $2$ euros;
        prêmios: uma cesta de $100$ euros e duas de $50$ euros.
        Calcule o ganho esperado de um bilhete e a margem da
        rifa; compare com a margem da casa na roleta europeia,
        $\frac{1}{37} \approx 2.7\,\%$. Por que a rifa sai
        impune?
  \item Uma seguradora mantém $1\,000$ apólices \omterm{def:g12:condprob:indep}{independentes}:
        \omterm{def:g10:proba:distribution}{probabilidade} de sinistro $0.01$ em cada uma, valor do
        sinistro $10\,000$ euros, prêmio $130$ euros. Calcule o
        lucro anual esperado \emph{e} o \omterm{def:g11:stat:variance}{desvio padrão} do total
        de sinistros. Compare os dois números: o que essa
        comparação obriga as seguradoras reais a manter?
  \item Finale --- a binomial como árbitro de decisões: a lista
        de verificação para reconhecê-la; a bússola
        $\E \pm \sigma$; as \omterm{def:g10:proba:distribution}{probabilidades} de cauda como preço
        do risco; e as duas ressalvas permanentes (a
        \omterm{def:g12:condprob:indep}{independência} é uma afirmação de modelagem, e são as
        caudas raras, não as \omterm{def:g11:stat:mean}{médias}, que causam a ruína). Uma
        frase para cada, com o ponteiro: a lei dos grandes
        números e a curva em sino, nos próximos capítulos,
        completam o regulamento do árbitro.
\end{enumerate}
\end{problem}
